\section{Method}
\label{sec:method}

Figure~\ref{fig:architecture} presents AgenticRL as a budgeted control process
over a fixed multimodal RAG environment. The retrieval modules, deployed vector
index, and generator are fixed during controller optimization. At each round,
a smaller trainable controller selects candidate evidence and chooses whether
to preserve the current evidence with \textsc{Accept} or redirect retrieval with
\textsc{Rewrite}. Its decisions are conditioned on a trajectory-aware state:
the latest non-empty retained set, which is the only memory the generator ever
sees, and an append-only history of failed retrievals, which only the controller
sees. A fixed three-round inference wrapper maintains both and exposes the
former to the generator once the budget is exhausted; the wrapper is
deterministic and holds no information of its own, and it exists because
retained evidence is useful only if a later round can reuse it.

\begin{figure*}[t]
  \centering
  \IfFileExists{figures/fig2_agenticrl_architecture.pdf}{%
    \includegraphics[width=0.96\textwidth]{fig2_agenticrl_architecture.pdf}}{%
    \framebox[0.95\textwidth]{\rule{0pt}{4.2cm}\ Figure~2 placeholder (system architecture)}}
  \caption{AgenticRL couples a frozen multimodal retrieval environment with a
  fixed-budget refinement wrapper. At each step, the controller makes
  candidate-level \textsc{Keep}/\textsc{Drop} decisions and a round-level
  \textsc{Accept}/\textsc{Rewrite} decision. The latest non-empty retained
  evidence $M_t^+$ is generation-visible, whereas append-only failed traces
  $M_t^-$ are controller-only. Training evaluates single-step action
  consequences and updates only the controller.}
  \label{fig:architecture}
\end{figure*}

\subsection{Problem Formulation and Fixed Environment}
\label{sec:formulation}

A multiple-choice instance is $x=(q,I_q,\mathcal{A})$, where $q$ is the
question, $I_q$ is the query image, and
$\mathcal{A}=\{a_1,\ldots,a_m\}$ is the option set. The controller may
condition on $\mathcal{A}$, but never observes the gold option label
$a^\star$. We distinguish the information available to each component:
\begin{equation}
\begin{aligned}
 x^{\mathrm{ctrl}}_t &=(q,I_q,\mathcal{A},z_t,\mathcal{C}_t,M_t^+,M_t^-),\\
 x^{\mathrm{ret}}_t  &=(z_t,I_q,\ell_x),\\
 x^{\mathrm{gen}}    &=(q,I_q,\mathcal{A},E),
\end{aligned}
\label{eq:component-inputs}
\end{equation}
where $z_t$ is the current retrieval query, $\ell_x$ is a test-time
anatomical routing signal, and $E$ is the final evidence context. In
particular, $M^-$ and $a^\star$ are never included in the generator input.

The environment contains text and image retrievers $R_T,R_I$, a deployed
vector index $\mathcal{I}$, and a vision-language generator $G$. Each indexed
unit stores text/image content, source and anatomical metadata, and
channel-specific scores. The gold answer is used only to compute utility after
executing a training action.

\subsection{Domain-Adaptive Multimodal Retrieval}
\label{sec:environment}

At round $t$, the text channel uses the current query $z_t$, whereas the image
channel uses the original query image $I_q$. In the optional modality-weighting
variant, a fixed routing function estimates their query-dependent
informativeness:
\begin{equation}
 (\alpha_t^T,\alpha_t^I)=f_{\mathrm{route}}(z_t,I_q),\qquad
 \alpha_t^T+\alpha_t^I=1.
 \label{eq:modality-weights}
\end{equation}
The routing function combines a text-density classifier, a visual-dependence
classifier, and a density ranker. Its parameters remain fixed, but its output is
recomputed after rewriting because $z_t$ changes. This module is evaluated in
separate front-end experiments; the reported 42.84\% memory-aware
configuration uses equal channel quotas without adaptive weighting.

The examination type and current query determine a gastrointestinal
anatomical partition $\ell_x$. The examination type is the endoscopy-type field
that the benchmark provides with each question as part of its input
specification; it is available at test time and is independent of the gold
answer, task label, and evaluation category. When it is absent or incompatible
with the current query, or when the examination covers no single
gastrointestinal site, no partition constraint is applied and retrieval
searches the full index. Otherwise it is applied as a Milvus metadata
pre-filter before approximate-nearest-neighbor search. The two channels
retrieve from the selected partition:
\begin{equation}
 \mathcal{C}_t^T=R_T(z_t;\mathcal{I}_{\ell_x}),\qquad
 \mathcal{C}_t^I=R_I(I_q;\mathcal{I}_{\ell_x}).
 \label{eq:dual-retrieval}
\end{equation}
When modality weighting is enabled, scores are normalized within each channel
before fusion,
\begin{equation}
 s_t(c)=\alpha_t^T\widetilde{s}_t^T(c)
       +\alpha_t^I\widetilde{s}_t^I(c).
 \label{eq:fusion}
\end{equation}
Each channel first retrieves 20 neighbors with \texttt{nprobe}=64. After
optional fusion, up to six candidates per channel form the 12-item controller
pool $\mathcal{C}_t$. Persistent evidence identifiers are used when
the same source is encountered again across rounds. Every rewrite traverses
this complete path again; no retriever, routing module, fusion parameter,
index entry, or generator parameter is updated.

\subsection{Budgeted Closed-Loop Controller}
\label{sec:controller}

The inference wrapper maintains states with different consumers. $M_t^+$
stores the current best evidence, defined as the latest non-empty
\textsc{Keep} set. It remains visible to the controller across rounds and is
the only memory supplied to the final generator. $M_t^-$ contains failed
retrieval traces, represented by the previous query and a deterministic compact
serialization of its rejected candidates. It is append-only, controller-visible,
and naturally bounded by the three-round budget. Thus $M^-$ is retrieval
history rather than negative answer evidence: it may guide reformulation but
cannot enter the answer context, alter index scores, or remove candidates.

With remaining decision budget $b_t$, the state is
\begin{equation}
 s_t=(x,z_t,\mathcal{C}_t,M_t^+,M_t^-,b_t).
 \label{eq:state}
\end{equation}
The structured action is
\begin{equation}
 a_t=(K_t,y_t,\widetilde z_{t+1}),\quad
 K_t\subseteq\mathcal{C}_t,\quad
 y_t\in\{\textsc{Accept},\textsc{Rewrite}\},
 \label{eq:action}
\end{equation}
where $\widetilde z_{t+1}$ is empty for \textsc{Accept} and contains the new
query for \textsc{Rewrite}.
The policy emits this action as JSON containing candidate decisions, the round
action, and an optional rewritten query. Candidate identifiers are range
checked and duplicates are removed by the shared parser.

Retained evidence follows replacement rather than union semantics:
\begin{equation}
 M_{t+1}^+=
 \begin{cases}
 K_t, & K_t\neq\varnothing,\\
 M_t^+, & K_t=\varnothing.
 \end{cases}
\label{eq:mplus}
\end{equation}
Consequently, a newly selected non-empty set becomes the current best state,
whereas an empty selection leaves the previous state unchanged. In the
memory-aware protocol, an initial \textsc{Rewrite} deterministically initializes
\begin{equation}
 M_1^+=\operatorname{Top3}(\mathcal{C}_0).
 \label{eq:mplus-init}
\end{equation}
This answer-independent initialization is reported explicitly and is not
treated as learned selection. Each rewrite also appends
\begin{equation}
 M_{t+1}^-=M_t^-\cup
 \{(z_t,\phi(\mathcal{C}_t\setminus K_t))\},
 \label{eq:mminus}
\end{equation}
where $\phi$ is the deterministic trace template used by the controller
prompt; it introduces no additional trainable summarizer.

If $y_t=\textsc{Accept}$, the wrapper updates $M_{t+1}^+$ and proceeds to the
next refinement step without another retrieval call. If
$y_t=\textsc{Rewrite}$, $z_{t+1}$ is executed through the same frozen retrieval
chain to obtain $\mathcal{C}_{t+1}$. One round denotes one controller decision,
and both branches consume one unit of the fixed three-round refinement budget.
The generator is called after the budget is exhausted using the final
$M_T^+$. A structured-output failure skips the corresponding state update and
continues to the next step. If no round produces retained evidence, the system
reuses the original single-round prediction. These wrapper rules are
deterministic and are not learned by the controller.
Thus, \textsc{Accept} means ``preserve evidence without retrieval at this
step,'' rather than early termination; termination is fixed by the refinement
budget. The learned round action controls retrieval allocation, not when the
complete system stops.

\subsection{Executed-Action Utility}
\label{sec:utility}

The controller is trained from evidence effects on the frozen generator. Let
$P_G(a^\star\mid q,I_q,E)$ be the next-token probability normalized over valid
option labels. Evidence utility is its gain over the same generator without
retrieval:
\begin{equation}
 U(E)=\log P_G(a^\star\mid q,I_q,E)
      -\log P_G(a^\star\mid q,I_q,\varnothing).
 \label{eq:utility}
\end{equation}
The subtraction provides a common diagnostic scale and cancels within a GRPO
group because all actions share the same no-evidence term.

For \textsc{Accept}, the executed context is the selected evidence
$E_t^A=K_t$. For \textsc{Rewrite}, the environment first executes the
new query and obtains $\mathcal{C}_{t+1}$. Its training-time context is the
five highest-ranked newly returned candidates:
\begin{equation}
 E_t^R=\operatorname{Top5}(\mathcal{C}_{t+1}),\qquad
 U_t^A=U(E_t^A),\quad U_t^R=U(E_t^R).
 \label{eq:branch-utility}
\end{equation}
This is a deterministic one-step retrieval-lookahead surrogate, not a
multi-round trajectory return: it evaluates whether a rewrite improves the
deployed retriever's immediate evidence. What a rewrite controls is the pool it
returns, whereas what a later round keeps from that pool is a selection decision
already supervised by $U_t^A$ on the same scale; moreover the retained state
after a rewrite is the top-ranked candidates of the new pool or a subset of
them, so the quality of that pool bounds what any later selection can retain.
Training states consequently do not unroll the three-round inference wrapper.
Although the two actions construct evidence differently, both are measured in
the same generator log-probability units.
This avoids a language-only rewrite score and penalizes fluent reformulations
that fail against the actual index.

\subsection{Three-Stage Controller Optimization}
\label{sec:optimization}

\paragraph{Supervised fine-tuning.}
Teacher actions initialize the structured schema and basic selection,
sufficiency, and rewrite behavior. Invalid JSON, out-of-range identifiers, and
schema-inconsistent rewrites are removed. SFT minimizes token-level negative
log-likelihood on controller outputs.

\paragraph{Rejection-sampling fine-tuning.}
Starting from the SFT policy, we sample structured actions and execute their
retrieval consequences. RFT retains policy actions with useful
generator-derived outcomes and balances rewrite-win states, evidence-selection
states, and trivial states. The retained actions are optimized with the
same supervised objective, moving the controller from teacher imitation toward
actions that work in the deployed environment.

\paragraph{Group-relative policy optimization.}
For each state, the policy samples a group of eight actions. After executing
each action, utilities $\{U_i\}_{i=1}^8$ produce normalized advantages
\begin{equation}
 A_i=\frac{U_i-\operatorname{mean}_j U_j}
 {\max(\operatorname{std}_j U_j,0.5)}.
 \label{eq:advantage}
\end{equation}
Groups with utility spread below $0.3$ are skipped. During training, two of the
eight samples are forced to explore \textsc{Rewrite}; these samples expose the
policy to retrieval consequences but their forced prefixes do not receive
policy-gradient credit. The remaining actions are sampled normally. The
clipped GRPO objective includes a KL penalty to the fixed RFT reference:
\begin{equation}
 \mathcal{L}_{\mathrm{GRPO}}=
 -\mathbb{E}_i\!\left[
 \min(r_iA_i,\operatorname{clip}(r_i,1-\epsilon,1+\epsilon)A_i)
 -\beta D_{\mathrm{KL}}(\pi_\theta\Vert\pi_{\mathrm{RFT}})
 \right].
 \label{eq:grpo}
\end{equation}
Only controller parameters are updated. In summary, SFT establishes valid
actions, RFT filters actions using executed retrieval outcomes, and GRPO
optimizes their relative utility without changing the retrieval environment or
generator.
